\section{Introduction}
\label{sec:intro}

Snow avalanches represent one of the most destructive natural hazards in alpine and high-mountain environments worldwide, posing severe threats to human life, critical transportation corridors, hydropower infrastructure, and mountain settlements across major mountain systems such as the European Alps, the Pamirs, the Himalayas, and Arctic mountain belts \citep{buehler2019automated, eckerstorfer2017remote, eaws2022size, techel2018spatial}. Rapid, systematic, and regional-scale mapping of avalanche release zones and debris deposits is vital for emergency search-and-rescue operations, hazard index updating, and retrospective validation of numerical avalanche dynamics models \citep{tompkin2021backscatter, muller2021synthetic}. However, conventional ground-based observations and optical satellite platforms (e.g., Sentinel-2, Landsat-8/9, PlanetScope) are frequently incapacitated during major avalanche cycles due to persistent cloud cover, winter blizzards, and polar night in high-latitude mountain regimes \citep{eckerstorfer2019near, vickers2016method}.

Synthetic Aperture Radar (SAR), particularly spaceborne C-band radar operating on the European Space Agency's (ESA) Sentinel-1 constellation, provides an indispensable all-weather, day-and-night observation system with a regular 6- to 12-day repeat cycle \citep{eckerstorfer2017remote, small2011flattening}. When snow avalanches release, dense slab flow and turbulent powder snow transport heavy snow masses down mountain chutes, depositing rough, granulated debris mounds in runout fans. This mechanical deposition alters the snowpack surface from smooth, undisturbed snow into rough, blocky debris, markedly increasing surface and volume backscatter ($\Delta \sigma^0 = \sigma^0_{\text{post}} - \sigma^0_{\text{pre}} > 0$) in both cross-polarization (VH) and co-polarization (VV) channels \citep{tompkin2021backscatter}.

Despite these promising physical principles, automated SAR change detection in rugged alpine terrain has historically suffered from acute false alarm rates and severe geometric distortions \citep{bianchi2021snow, gatti2026avalanche}. In steep topography, the side-looking imaging geometry of SAR creates complex radiometric distortions:
\begin{enumerate}
    \item \textbf{Local Incidence Angle (LIA) Compression:} Mountain slopes tilted directly toward the radar antenna experience extreme ground-range compression (foreshortening) and layover, resulting in artificially high backscatter values that mimic avalanche debris. Conversely, slopes tilted away from the sensor produce high incidence angles and radar shadow, where signal-to-noise ratios collapse \citep{small2011flattening, tompkin2021backscatter}.
    \item \textbf{Directional Slope-Aspect Alignment:} Sentinel-1 operates in near-polar orbits with fixed ascending (look azimuth $\phi_{\text{look}} \approx 78^\circ$) and descending ($\phi_{\text{look}} \approx 282^\circ$) lines of sight. As established by \citet{tompkin2021backscatter}, avalanche backscatter contrast is highly anisotropic with respect to the angle between slope aspect and radar look direction: slopes facing away from the radar exhibit the highest relative brightness anomalies ($55^\circ \pm 20^\circ$ LIA), while foreslopes generate severe topographic clutter. In complex alpine relief, steep topography also influences winter terrain roughness and snow drift accumulation \citep{veitinger2014large}.
    \item \textbf{Slope-Prior Sensitivity:} Flat valley bottoms ($<5^\circ$) and very steep rock faces ($>65^\circ$) can contain clutter, but they also contain genuine mapped debris through runout, terrain-model error, or geolocation effects. Any slope prior must therefore be treated as a heuristic rather than a physical-validity mask.
\end{enumerate}

Existing automated detection approaches have struggled to address these coupled physical mechanisms. Early methods relied on empirical thresholding of backscatter differences ($\Delta \sigma^0$) combined with heuristic slope filters \citep{eckerstorfer2019near, vickers2016method}; \citet{liu2021object} subsequently combined object-based Sentinel-1 classification with topographic causative indices using SVM and logistic-regression mapping. Recent deep-learning work introduced Siamese convolutional networks, attention-augmented change detection \citep{daudt2018fully, zhang2018road, chen2021remote, bandara2022transformer, gatti2026avalanche}, and promptable foundation models \citep{gelato2026promptable}. \citet{bianchi2021snow} investigated convolutional segmentation with slope and potential-angle-of-reach masks, while \citet{gatti2026avalanche} explored multimodal Swin Transformer concatenation on AvalCD. TopoRadar-Net differs from the object-based SVM/LR workflow of \citet{liu2021object} through learned SAR-query/topography-key-value fusion and differs from single-configuration studies through cross-system, multi-seed stability and negative-result characterization. In the wider natural-hazards literature, geomorphological assessment and deep-learning susceptibility studies show severe transfer hurdles across catchments \citep{wang2023evaluation, sajjad2024assessing, zhou2023deformation, guo2025forecasting, guo2024landslide, garova2025general, liu2026integrating}. Our controlled comparisons evaluate whether auxiliary radar-topographic cross-attention and continuous aspect-look coordinates provide measurable advantages over standard architectures.

To overcome these fundamental limitations, we propose \textbf{TopoRadar-Net}, a lightweight, physics-guided cross-attention architecture specifically engineered for alpine SAR change detection. Rather than relying on black-box concatenation, TopoRadar-Net introduces four methodological, theoretical, and operational innovations:
\begin{itemize}
    \item \textbf{Theoretical Formulation of 2D Convolutional Inadequacy:} We formulate the geometric limitations of shift-invariant 2D linear convolutions operating purely on SAR intensity space under the non-stationary radiometric projection field $\frac{\cos \psi}{\sin \theta_{\text{inc}}}$. We show that projecting terrain into a continuous 4D aspect-look manifold $\mathcal{M}$ provides continuous orientation coordinates to condition feature representations without circular angular branch-cut discontinuities.
    \item \textbf{Multi-Scale Radar-Topographic Cross-Attention Block (RTCAB):} We design a cross-attention mechanism wherein multi-scale bitemporal SAR change features serve as queries to interrogate regionally-pooled topographic geometry keys and values, conditioning spatial feature representations on terrain context.
    \item \textbf{Heuristic Geomorphic Slope Prior (Geo-Loss):} We evaluate a soft penalty on predictions below $5^\circ$ and above $65^\circ$. A construct audit shows that the mask also covers genuine positive labels, so its ablation is interpreted as component sensitivity and its subgroup trade-off is reported explicitly.
    \item \textbf{Cross-System Generalization and Conceptual Decision Support:} We evaluate TopoRadar-Net across four mountain systems comprising 7 Sentinel-1 scenes and 993 avalanche polygons. Across three seeds, the pooled margin over Attention U-Net is $+0.38\%$ ($p=0.869$), while the margin over SiamUNet-diff is significant ($+8.92\%$, $p=0.021$). Exploratory Seed-42 backslope differences do not persist as stable three-seed effects (Troms{\o}: $-0.34\% \pm 6.28\%$; Pamir: $+1.20\% \pm 7.73\%$), delimiting the evidence for directional conditioning. An explicitly unvalidated triage framework accompanies a $30.0\%$ Pamir false-alarm difference ($72.9\text{ ha}$ nominal-grid equivalent; $55.6\text{ ha}$ affine physical difference) and $1.80\text{ s}$ model-forward latency.
\end{itemize}

The remainder of this article is organized as follows: Section~\ref{sec:data} describes the regional settings, the AvalCD benchmark dataset, and data preprocessing. Section~\ref{sec:method} details the radar-topographic physics, the TopoRadar-Net architecture, and the loss formulation. Section~\ref{sec:results} presents comparative baseline performance, cross-region zero-shot evaluations, instance-level EAWS detection, and ablation studies. Section~\ref{sec:discussion} discusses physical mechanisms, operational mountain safety implications, and study limitations. Section~\ref{sec:conclusion} concludes the paper.
